\chapter{A global controlled surgery flow}
\label{ch:global-surgery-v4}

This chapter constructs the changing-carrier Ricci flow used by the later
extinction argument.  The word ``global'' refers to the time domain
$[0,\infty)$, not to a flow on one fixed manifold: at a surgery time the
carrier is cut and capped, and all identifications after that time are part
of the data.  The construction has two logically different outputs.  First,
an already controlled finite flow can be continued on its own carrier
history.  Second, a normalized initial metric produces one complete flow
with the same parameter schedule.  Keeping these branches separate is what
prevents the global theorem from silently replacing a given solution.

Throughout, $M$ is a compact smooth three-manifold and $g_0$ is the normalized
initial metric selected in the preceding chapters.  A raw flow is a
changing-carrier object

\[
        F=(M_t,g(t),\mathcal J_F,\mathcal E_F),
\]

where $\mathcal J_F\subset\mathbb R$ is its time domain,
$\mathcal E_F$ is the event family, and $M_t$ is allowed to change at an
event.  We write $J_F$ for the set of surgery times.  On a nonsurgery slab
$g(t)$ solves $\partial_tg=-2\operatorname{Ric}$; an event carries its
incoming slice, the retained region, the cap charts, the post-surgery slice,
and the comparison maps.  Admissibility and pinching have the meanings in
Section~\ref{sec:surgery-carriers}. Canonical control at radius $r(t)$
requires the strong canonical alternatives at every point with
$R\geq r(t)^{-2}$. Noncollapsing is imposed on surviving backward
parabolic balls of radius at most $\epsilon$, with the positive-component
exception described there.

We also impose the terminal event condition produced by
Theorem~\ref{thm:branch-continuation}. Put $\rho=\delta(T)r(T)$ and
let $\Omega_{\rm core}$ be the union of components of the regular terminal
region meeting $\{R\leq\rho^{-2}\}$. At a nonempty event the retained
region, transported to that limit, is exactly $\Omega_{\rm core}$ minus
the chosen escaping neck ends. Each end is the complementary component
beyond its central sphere, contains the positive half-neck, misses the
negative half-neck, and escapes every compact set. At a vanishing event,
for each point $x$ of the last ordinary slab there are $L_x>\rho^{-2}$
and $s_x<T$ such that $R(x,t)>L_x$ for $s_x\leq t<T$.
This condition is asserted only at events in the interval being observed.

\section{The order of choices}
\label{sec:global-choice-order}

The numerical order is part of the theorem.  It is useful to state it before
any induction.  A metric-surgery constant package $K$ contains the local
cutoff $K.\delta_0$, curvature constant $K.C_0$, and the scale $K.R_0$.
A setup for $K$ consists of a standard cap flow, numbers
$\epsilon,C$, and a common selector $h(\rho,\delta)$.  It satisfies

\[
  0<\epsilon<\tfrac12,\qquad
  \epsilon\leq\min\{1/200,K.\delta_0/2\},\qquad
  0<C_0\leq C,\qquad 1\leq C,
\]

and, for positive $\rho,\delta$,
$0<h(\rho,\delta)\leq\min\{K.\delta_0\rho,\rho\delta\}$.  The selector is monotone in both
arguments.  The surgery height at time $t$ is always

\begin{equation}
  h_t=h(\delta(t)r(t),\delta(t)).
  \label{eq:global-height}
\end{equation}

Fix the standard initial metric and the particular standard flow and
persistence theorem of Section~\ref{sec:standard-cap}. Calibration
chooses, in this order, the small-neck threshold, the
bounded-distance and dense-time thresholds, the neck-and-cap topology theorem,
the common singular-limit accuracy, the gluing parameter $\beta$, the
canonical constant

\[
                         C=\max\{C_{\kappa},C_{\rm standard}+1\},
\]

the model four-jet estimates, the initial noncollapsing constant $\kappa_0$,
and finally the deep-horn height selector of
Corollary~\ref{cor:singular-selector}. Here $C_\kappa$ controls the ancient
canonical alternatives at accuracies $\epsilon$ and $2\epsilon$, with
their projective-plane-product exception; $C_{\rm standard}$ controls
the chosen standard flow at accuracy $\beta\epsilon/3$.
Four intrinsic spatial jets of a neck or round model bound
$|\nabla R|/R^{3/2}$ and $|\Delta R+2|\operatorname{Ric}|^2|/R^2$.
They do not assert a bound across a surgery jump.
The selector is chosen only after this analytic coefficient is fixed.

The number $0<\beta<1/2$ controls gluing of two evolving necks. The
recent neck is $\beta\epsilon$-close to the shrinking cylinder; the
older strong neck has accuracy $\beta\epsilon/2$, normalized by its own
joining-time scalar. An exact joining isometry identifies their metrics
and centers. Their common recent spatial chart, restricted to length
$\epsilon^{-1}$, then gives an $\epsilon$-neck on normalized time
$(-1,0]$, smooth through the joining slice
\cite[Proposition 15.2, pp.~353--354]{MT2007}.

Here is the construction with its time and scale matching. Normalize the
final central scalar to one, write $s=-d$ for the joining time, and let
$\chi$ be the recent spatial chart and $I$ the joining isometry. The
joined tensor is $\chi^*g_{\rm recent}(t)$ for $t\geq s$ and
$(I\circ\chi)^*g_{\rm older}(t)$ for $t<s$. If $d\geq1$, the
recent piece already covers the required interval. If $0<d<1$, sufficiently
small $\beta\epsilon$ controls four spatial jets and gives, at the
central point throughout $[s,0]$, $R>1/4$ and $|\Delta R|<1/100$.
Since $|\operatorname{Ric}|^2\geq R^2/3$, the scalar equation makes
$R$ strictly increasing there. Thus $1/4<R(s)<R(0)=1$.
The older neck's squared scale is exactly $\lambda=R(s)^{-1}$, so
$1<\lambda<4$; its strong-neck interval reaches back to $s-\lambda<-1$.
Its normalized time corresponding to physical time $t$ is
$\tau=(t-s)/\lambda$. The coordinate transition $\phi$ from the
recent chart through $I$ to the older chart consequently satisfies
$A=\lambda\phi^*C(0)$, where $A=\chi^*g_{\rm recent}(s)$ and
$C(\tau)$ is the older normalized tensor.

For completeness, the accuracy is uniform over all such pairs. Otherwise
choose failing pairs with $\beta\to0$ and points and times at which the
required squared jet error is at least $\epsilon^2/2$. Pass to a
subsequence with $d\to d_\infty\in[0,1]$; scalar readout from the
joining metric gives $\lambda\to1+d_\infty$. In coordinates centered
at each tested point, model ellipticity and the exact metric identity
bound $D\phi$. Differentiating the Christoffel transformation formula
then bounds every fixed higher derivative of $\phi$ by induction.
These bounds require only the actual open coordinate neighborhoods at
the tested points. Write the shrinking-cylinder tensor as
$B(t)=B(0)+tS$, with $S=-2\operatorname{Ric}(B(t))$ independent of
$t$. Ricci naturality at the joining slice and convergence of the
metric jets imply $\phi^*S-S\to0$ in every fixed jet order. The exact
identity
\begin{align*}
 \lambda\phi^*C(\tau)-B(t)
 ={}&A-B(s)-\lambda\phi^*(C(0)-B(0))\\
 &+\lambda\phi^*(C(\tau)-B(\tau))
 +(t-s)(\phi^*S-S)
\end{align*}
now makes the older-piece error tend to zero through the fixed order
$\lfloor\epsilon^{-1}\rfloor$. The recent-piece error already tends
to zero. This contradicts the chosen violating samples and fixes
$\beta$ before the input pair.

Finally, equality of the joining metrics on the open chart matches all
their spatial derivatives at $s$. On each side the coefficient equation
$\partial_tg=-2\operatorname{Ric}(g)$ expresses time derivatives in
terms of spatial derivatives and the inverse metric. Induction therefore
matches all mixed derivatives across $s$, and pasting the two smooth
families gives joint smoothness there. Restrict $\chi$ to axial interval
$(-\epsilon^{-1},\epsilon^{-1})$; its coordinate map and inverse are
unchanged, and the resulting tensor has the required comparison on all
of $(-1,0]$.
The initial cutoff is

\begin{equation}
 \Delta_0=\min\left\{\frac{\beta\epsilon}{3},\delta_{13},
       K_0^{-1},D_0^{-1}\right\},
 \label{eq:global-delta0}
\end{equation}

where $\delta_{13}>0$ is an applicability cutoff for metric surgery,
$K_0$ bounds the volume of the standard metric's closed core, and
$D_0^{-1}\leq R(g_0^{\rm cap})\leq D_0$. All four entries are positive.
The normalized initial-flow estimate supplies a Ricci flow on
$[0,1/16]$ with curvature bounded by $2$ and volume lower bound
$\operatorname{Vol}(B(x,r))\geq\kappa_0r^3$ for $r\leq\epsilon$; ordinary
uniqueness and maximality identify it with the eventual flow and exclude an
early surgery \cite[Claim 15.1, p.~353; Definition 15.7, p.~360]{MT2007}.

The gluing assertion used here fixes $\epsilon$ before $\beta$. It does
not require one $\beta$ to work uniformly for every accuracy. Subsequent
cutoff restrictions preserve this same fixed-accuracy setup.

\section{Prefixes and the global numerical schedule}
\label{sec:global-prefixes}

For a positive integer $i$, a parameter prefix $p$ consists of arrays

\[
 (r_0,\ldots,r_i),\qquad (\kappa_0,\ldots,\kappa_i),\qquad
 (\Delta_0,\ldots,\Delta_i),
\]

all positive and antitone, with $r_0=\epsilon$, $r_j\leq\epsilon$, and
$\Delta_j\leq K.\delta_0$.  The $j$th entry controls the half-open interval

\begin{equation}
 E_j=\begin{cases}
 [0,2^0/32),&j=0,\\
 [2^{j-1}/32,2^j/32),&j>0.
 \end{cases}
 \label{eq:global-entry}
\end{equation}

The endpoint $T_j=2^j/32$ is the start of the next physical epoch.  A
prefix is seed-compatible when it uses the fixed setup,
$\kappa_0$ is its initial noncollapsing constant, and $\Delta_0$ is the value in
\eqref{eq:global-delta0}.

The noncollapsing induction has the following quantified form:

\begin{quote}
for every seed-compatible $p$ there is a positive $\kappa_{\rm new}$ and a
function $\operatorname{cutoff}(r)$ such that
\begin{align*}
0<\kappa_{\rm new}&\leq\kappa_i,\\
0<\operatorname{cutoff}(r)&\leq\Delta_i
\end{align*}
for $0<r\leq r_i$, and every observed next-epoch flow satisfying the old
prefix controls, the terminal policy, the canonical bound at $r$, and
$\delta\leq\operatorname{cutoff}(r)$ is $\kappa_{\rm new}$-noncollapsed.
\end{quote}

Here ``noncollapsed'' is the actual cylinder statement: if a point is not in
a positive component, every surviving backward cylinder of radius
$0<s\leq\epsilon$ on $[t-s^2,t]$, based on the whole terminal ball
$B_{g(t)}(x,s)$ and equal to the identity there at $t$, with
$|\mathrm{Rm}|\leq s^{-2}$ has

\begin{equation}
 \operatorname{Vol}_{g(t)}B(x,s)\geq \kappa_{\rm new}s^3.
 \label{eq:global-noncollapse}
\end{equation}

A positive-action barrier first keeps a minimizing
backward path out of a cap window.  A finite-gauge direct method then attains
the action infimum on a compact confined region.  Free-endpoint variation and
the capped slice-value comparison produce a minimizing region.  The old seed
cylinder is joined to it by the comparison paths of Claim 16.27; the stable
image has positive measure. Theorem~\ref{thm:m15-noncollapse} gives the volume estimate at the
chosen radius.  The half-radius passage loses the explicit factor $1/8$.
Low scalar points and radii below the auxiliary scale are treated by a larger
controlled cylinder and Bishop--Gromov comparison.  Positive components are
excluded throughout the same history, which is precisely Remark 16.2's
exception.

The order of the auxiliary choices is part of this argument.  For a proposed
next radius $r_{\rm next}$, the producer first fixes
$B=\max\{1,B_{\rm seed}\}$, where $B_{\rm seed}$ is the fixed analytic
coefficient for the initial and standard-flow estimates, and the positive
action budget, then takes the
test radius $\rho=\rho(B,r_{\rm next})\leq r_{\rm next}$.  It next chooses
the action-barrier data: a cap radius $A$, comparison tolerance $\eta$,
normalized time $1/2<\theta<1$, and scalar lower coefficient $c>0$.
To see their order, write $L_+$ for the action budget after adding the
scalar-floor correction. Choose a safe duration $a>0$ with
$a\leq\rho^2/12$ and $L_+<\rho^2/(16\sqrt a)$. The standard flow's
scalar lower rate supplies $c$; choose $\theta$ so that
\[
 \frac{L_+}{\sqrt a}
 <-\frac c2\bigl(\log(1-\theta)+\log2\bigr).
\]
On $[0,\theta]$ the model metric is at least $\mu>0$ times its birth
metric. Choose $A$ to contain the core buffer and satisfy
$L_+/\sqrt a<\mu A^2/4$. Finally choose $\eta$ so the actual cap
comparison retains these metric and scalar bounds. Escape through the
side then costs too much kinetic action, and reaching the top costs too
much scalar action. Only after these choices select the common
cap cutoff $\delta_0$ from the seed, scalar and cap inequalities.  The actual
cutoff is

\[
 \delta=\min\{\delta_0,\,\delta_{\rm cap}(\rho/2)\}>0.
\]

The confinement statement concerns paths below the chosen action barrier.
Comparison paths to the closure of an open neighborhood $V_0$ in the old seed
slice have action at most half that barrier.  An infimizing sequence can
therefore be restricted to this sublevel; leaving the compact spacetime cage
or meeting a forbidden cap window would cost more than its action bound.
In finitely many regular gauges, the square-root-time energy estimate gives
weak compactness of velocities and uniform compactness of positions.  The
kinetic integral is lower semicontinuous and the scalar term converges, so
the limit attains the infimum.  Endpoint variations and the capped
slice-value comparison supply the minimizing region.  Every minimizer to
the closure of $V_0$ has the same action bound, because it has no greater
action than the explicit comparison path.

Here is the additional step needed to apply generalized noncollapsing.
Fix the exponential family $E$ supplied by this minimizing region and the
backward duration $\tau=b^2$, where $b=\sqrt\tau\geq\epsilon$. The seed
construction keeps the endpoint time in the interior of the regular history
and bounds this duration by the elapsed time since the preceding epoch began.
All minimizing initial vectors with
endpoint in $\overline V_0$ lie in a compact set.  Indeed, the curvature and
scalar-gradient bounds in the cage bound the square-root-time velocities;
at the initial endpoint their squared norm is $4|Z|^2$.  Limits of these
vectors continue through the same compact phase region.  Recovery of nearby
endpoints by local comparison paths proves that the limiting branch still
minimizes.  Thus the minimizing vectors form a closed bounded set in the
finite-dimensional initial tangent space, and every minimizer is represented
in this same $E$.

Outside the critical image of the full surviving exponential map, each
minimizing vector has a smooth inverse branch.  Compactness supplies
finitely many such branches which capture all minimizers at nearby
endpoints: otherwise escaping minimizers would have a convergent subsequence
over the original endpoint outside their union.  The minimum action nearby
is the minimum of the finitely many smooth branch actions, hence is locally
Lipschitz.  The Sard and Rademacher argument for the stable set now applies:
at a differentiability point all minimizing branch actions have the same
derivative, hence the same terminal momentum. Backward uniqueness for their
ordinary differential equation makes the minimizing branches equal.
Compact capture then confines all nearby minimizers to this one inverse
branch, giving a stable neighborhood. Thus outside a null set the minimizer
is unique, noncritical and stable. In particular
the stable image covers $V_0$ up to a null set.  This conclusion uses the
compact capture just proved, and is not a general surjectivity assertion
about a generalized exponential map.

Let $H$ be the stable source and put $W=H\cap E_\tau^{-1}(V_0)$.  It is open,
and $E_\tau(W)=V_0\cap E_\tau(H)$, so its image has the full volume of $V_0$.
If $L$ is the action budget and $s_*$ is the old seed-image radius, the
construction supplies the uniform bounds
\[
 l|_{E_\tau(W)}\leq l_0=\frac{L}{4\epsilon},\qquad
 \operatorname{Vol}(E_\tau(W))\geq
 V=\frac{\kappa_{\rm old}s_*^3}{8},\qquad
 \max\{\epsilon^2,(s/2)^2\}\leq\tau\leq\overline\tau,
\]
where $s$ is the radius of the ball whose noncollapsing is being tested.
Together with the compact terminal ball and the actual curvature-controlled
half-radius cylinder, these are precisely the hypotheses of generalized
noncollapsing.  The constants $l_0,V,\overline\tau$ depend on the old prefix,
before the tested point and its flow are chosen.  Applying that theorem and
including the half-radius factor gives the asserted new noncollapsing
constant.  The action and stable-image arguments correspond to
\cite[Corollary 6.67 and Claims 6.68--6.69, pp.~139--140;
Proposition 16.1 and Lemma 16.15, pp.~367--368, 379--382, 394]{MT2007}.

Canonical induction uses this particular noncollapsing extension at $p$.
It supplies $r_{\rm next}$ and
$\Delta_{\rm next}$ with

\[
  0<r_{\rm next}\leq r_i,\qquad
  0<\Delta_{\rm next}\leq\operatorname{cutoff}(r_{\rm next}),
\]

and proves the canonical-neighborhood predicate on the next observed epoch.
The canonical conclusion holds for the same cutoff and noncollapsing
constant just selected.

We prove the canonical assertion by contradiction. Fix the old prefix and
its selected noncollapsing extension. If no next radius works, choose
$r_n\downarrow0$ and cutoffs $\delta_n\downarrow0$, with the cutoffs
chosen before their counterexample flows. At each stage also impose the
first $n$ cap-persistence bounds on duration, scalar curvature and distance.
Finite minima suffice for this choice. These bounds will survive every
subsequence, including the one selected later for a geometric limit.

In each counterexample take the infimum of the canonical failure times.
The standard-cap cover closes the failure at that infimum, so there is an
actual bad point $(x_n,t_n)$ with scalar $Q_n\geq r_n^{-2}$ and canonical
control at all strictly earlier times. The threshold is an inequality:
we normalize by the actual $Q_n$. Since $r_n\leq1/200$, the closed neck
window of length $Q_n^{-1}\leq r_n^2$ lies in the old overlap interval.

The volume input here must include positive components, which the preceding
noncollapsing assertion excluded. The residual failures lie in small
high-curvature whole positive components. Trace such a component backward
to its birth and to the onset of its high-curvature regime. If that onset
is at most $r_n^2/(32A_a)$ old, an accessible seed set has volume at least
$V\,a^{3/2}$ at age $a$ and reduced length at most $2$. Generalized
noncollapsing on that actual cylinder gives density
$\kappa_{\rm seed}/64$. If it is older, a
retained whole-component cylinder extends for more than three safe time
steps. Confinement and the action barrier apply there; distinguishing an
old birth from a recent birth gives density
$\min\{\kappa_{\rm old},\kappa_{\rm seed}/8\}$. Here $A_a$ is the
fixed analytic coefficient. Take the minimum of these positive densities
before choosing $r_n$, then use auxiliary radius $r_n/(8A_a)$ and reduce
the cutoff. Combining these estimates with the old prefix gives one volume
constant on the new epoch and its preceding physical buffer of length
$1/128$, including its endpoint. Thus the selected bad points have full
volume control, not merely control away from positive components.

Rescale each actual history by $Q_n$ about $(x_n,t_n)$. Pinching makes
negative curvature disappear, while the strict-past canonical estimates,
analytic bounds and buffered volume give compactness on controlled backward
cylinders. A maximal-cylinder search can stop at a cap at either positive
or zero backward duration. The preselected cap estimates make the original
center canonical in both cases, contrary to its selection. Consequently the
search supplies regular cylinders. The retained-terminal curvature argument
gives a scalar ceiling $B$ before the spatial radius, and analytic control
then gives a common backward time $\tau>0$. Taking a test radius
$\rho\leq R$, $\rho^2\leq\tau$, $B\rho^2\leq1$ converts buffered
volume to the base-ball lower bound for compactness. Choose a maximal common backward
horizon $H\in(0,\infty]$ for a subsequence, retaining the actual source
maps and all the finite cap bounds.

If $H=\infty$, compactness gives a complete nonflat ancient limit with
nonnegative curvature. Its noncollapse holds at every scale. Indeed, for a
limit ball of radius $r$ and $0<\theta<1$, use source curvature radius
$\theta^2r$, source volume radius $\theta r$ and an intermediate radius
$(1+\theta)r/2$. Compact containment and convergence give lower volume
$\kappa\theta^9r^3$; letting $\theta\uparrow1$ gives $\kappa r^3$.
The ancient classification therefore supplies a neck, cap, compact positive
component or round alternative. Capture its whole required spatial and
backward-time domain before transferring it to the same original bad point.
The accuracy margins make that point canonical, a contradiction.

Suppose instead that $H<\infty$. The finite-endpoint construction of
Section~\ref{sec:singular-blowup} first gives compatible local backward
extensions on an exhaustion and a smooth complete metric at time $-H$.
It dominates the complete time-zero metric. A single further subsequence
retains the endpoint comparison maps. At high-scalar endpoint points,
strict-past canonical neighborhoods and convergence of their jets give
related necks unless the component is compact. Strictly positive curvature
then gives a global scalar bound by the small-neck scale estimate. In the
null-plane case, the orientation cover and parallel null directions from
the local backward germs, together with a related neck, bound the transverse
curvature. The identically scalar-flat case is flat. Thus the endpoint has
a global bound $B_*$ before a common extension interval is chosen.

Put $L=2\max\{1,3B_*\}$ and choose
\[
 d=\min\{H/4,1/(256A_aL)\}>0.
\]
This choice precedes every spatial radius and pinching error. Old controls
at horizon $H-d$ give $B_{\rm old}$; use the common bound
$B_+=\max\{13\max\{4L,1\},B_{\rm old}\}$ on the proposed longer
window. Search backward on the original histories while preserving their
old comparison maps. For a terminal ball of radius $R$, a cap contact has
distance budget
\[
 D=2\sqrt{2\exp(6B_+(H+4d))}\,R.
\]
One of the finite cap bounds imposed before selecting the sequence dominates
this distance, the duration and the curvature bound. Such contact would
again make the original center canonical, so it cannot occur. The same
physical sequence consequently carries controlled cylinders of every radius
through time $-H-d$, contradicting maximality of $H$. This proves the
canonical extension, with its radius chosen after the uniform volume
constant and its cutoff chosen before the observed flow
\cite[Proposition 17.1, pp.~395--409]{MT2007}.

Starting with the seed prefix $p_0$, construct $p_{n+1}$ by adjoining
the radius and constants supplied by the noncollapsing and then canonical
induction at $p_n$. The prefix $p_n$ has last index $n+1$.
The old $r$ and $\kappa$ entries are copied exactly. The old last
$\Delta$ entry is replaced by the selected next cutoff, and this same
cutoff is appended as the new last entry. Thus, for every old index $j$,

\begin{equation}
 r_j^{(n+1)}=r_j^{(n)},\qquad
 \kappa_j^{(n+1)}=\kappa_j^{(n)},\qquad
\Delta_j^{(n+1)}\leq\Delta_j^{(n)}.
 \label{eq:global-stability}
\end{equation}

For $j\leq n$ the last inequality is an equality: only the old last
entry can change. Thus each coordinate stabilizes after its one tightening.
The schedule reads coordinate $j$ from stage $j$ after that tightening:
$r(j)=r_j^{(j)}$, $\kappa(j)=\kappa_j^{(j)}$, and
$\Delta(j)=\Delta_j^{(j)}$.  The stability equations imply that these are
positive antitone sequences and that every finite restriction agrees with the
corresponding prefix: for $j\leq n+1$ there is equality for $r,\kappa$
and $\Delta(j)\leq\Delta_j^{(n)}$, with equality also for $\Delta$ when
$j\leq n$. Positivity is inherited from that finite stage; no infinite
infimum of positive cutoffs is taken. The overlap height inequality follows from
monotonicity of $h$ and the three antitone sequences.  This is the concrete
meaning of ``later refinements preserve earlier controls''.

For a positive antitone control function $\delta:[0,\infty)\to(0,\infty)$,
let $j(t)$ be the unique index with $t\in E_{j(t)}$ and define

\begin{align}
 r(t)&=r(j(t)),\\
 \kappa(t)&=\kappa(j(t)),\\
 h(t)&=h(\delta(t)r(t),\delta(t)).
 \end{align}

On $E_j$ this gives $r(t)=r(j)$ and $\kappa(t)=\kappa(j)$.
If $\delta(t)\leq\Delta(j)$ there, then the full parameter record agrees
with the schedule used in continuation and global assembly:

\begin{equation}
 (r_F(t),\kappa_F(t),\delta_F(t),h_F(t))
 =\bigl(r(j),\kappa(j),\delta(t),h(\delta(t)r(j),\delta(t))\bigr).
 \label{eq:global-profile}
\end{equation}

The half-open convention is essential at an epoch boundary: the old value is
used strictly before the boundary, while the next value is used at and
after it.

\section{One singular epoch and the next frontier}
\label{sec:global-epoch}

Fix a compatible prefix $p$ with index $i>0$.  Suppose $F$ has time domain
$[0,H)$ with

\[
 T_i\leq H<T_{i+1},
\]

and a preterminal ordinary slab whose curvature is unbounded as $t\nearrow H$.
The observation $O$ has horizon $H$, old-prefix controls, the terminal event
policy, the next-epoch profile identities, and the overlap inequality
$\delta\leq\Delta_{\rm next}$. The one-epoch continuation gives an extension
$E$ and a new observation $O'$ with

\begin{equation}
 H<H'\leq T_{i+1},\qquad H\in J_E,\qquad
 J_E\cap(H,H')=\varnothing.
 \label{eq:global-progress}
\end{equation}

If $H'<T_{i+1}$, then $\mathcal J_E=[0,H')$ and a new preterminal slab
starts at $H$.  If $H'=T_{i+1}$, the returned flow is not truncated merely
because it reaches the epoch boundary; should its domain end there, the slab
starting at $H$ is still recorded.  The extension retains the old prefix
controls, is admissible and pinched, carries the selected canonical and
noncollapsing predicates on $[0,H')$, and keeps the standard cap model tied to
the actual observation.  The terminal policy for the new event is supplied on
the singleton $\{H\}$ and is combined with the old policy; no policy for
unobserved future events is inferred.

Here is the substantive bridge behind this statement.  The retained regular
history identifies the last ordinary slab with the generalized spacetime used
by the singular-limit theorem. Theorem~\ref{thm:singular-m31} supplies
a regular reference and a limit metric; Corollary~\ref{cor:singular-selector}
supplies a deep-horn scale at the already fixed coefficient. The
terminal core is the union of components meeting the terminal scalar
sublevel set at the selected scale. The bridge map is the identity on the old carrier after the explicit
reference equality, so the limit identity is an equality of the selected
limit and its comparison maps. Apply Theorem~\ref{thm:metric-surgery}
with the flow's actual constants and standard model.
Theorem~\ref{thm:branch-continuation} constructs the surgery and its maximal
ordinary restart. The frontier
lemma observes the endpoint of that restart or the epoch boundary, preserving
the same standard flow and the preterminal slab.  Parameter equality on the
extension transports the old profile; the canonical induction first gives
the canonical bound, and the noncollapsing induction applies to that same
branch with that bound.

The scalar derivative estimates used in constructing the singular-limit input live at the
auxiliary history radius and coefficient.  They are not an extra hypothesis
on the public epoch-extension statement and do not assert a derivative bound
through a surgery jump.

\section{Volume loss at the corrected scale}
\label{sec:global-volume-loss}

The event-count argument needs a quantitative loss at every genuine cap
event.  For a nonempty post-surgery slice at $T$, let $e_T$ be its event,
$t_T^-$ its left reference time, $A_T^-$ and $A_T^+$ the retained pre- and
post-regions, $g_T^*$ the regular-limit metric, and $R_{T,i}$ the specified
positive half of the $i$th neck. Let $q_T$ identify the retained incoming
region with its image in the regular limit. Write
$v_F(t)=\operatorname{Vol}_{g(t)}(M_t)$.  The volume-loss certificate
contains a loss $\ell_T\in[0,\infty]$ and a
finite left limit $L_T$ satisfying

\begin{align}
 v_F(t)&\longrightarrow L_T
   &&(t\nearrow T),\\
 \operatorname{Vol}_{g_T^*}(M_T^*)&\leq L_T,\\
 \operatorname{Vol}_{g(T)}(A_T^+)&=
   \operatorname{Vol}_{g_T^*}(q_T(A_T^-)),\\
 v_F(T)+\ell_T&\leq \operatorname{Vol}_{g_T^*}(M_T^*).
 \label{eq:global-event-volume}
\end{align}

The event alternative is the geometric source of positivity.  Either some
cap index $i$ satisfies $\ell_T>0$ and

\begin{equation}
 \operatorname{Vol}_{g(T)}(C_{T,i})+\ell_T
 \leq \operatorname{Vol}_{g_T^*}(R_{T,i}),
 \label{eq:global-neck-loss}
\end{equation}

or the event has zero caps, $\ell_T=0$, and a connected component outside
$A_T^-$ has strictly positive incoming volume.  Vanishing events have empty
post-volume and are bounded by their supplied left-limit volume.  The
zero-cap alternative is deliberately not forced into a positive loss estimate.

The corrected estimate is cubic:

\begin{equation}
 \ell_T\geq c\,\frac{h(T)^3}{\delta(T)}
 \qquad\text{when the event has at least one cap}.
 \label{eq:global-cubic-loss}
\end{equation}

For each finite horizon $H$, there is also a common $\sigma_H>0$ with
$\sigma_H\leq\ell_T$ for all capped events in
$[0,H]$.  To prove this, write $n_T$ for the number of inserted caps.
In the normalized neck chart the positive axial interval has length
$\delta(T)^{-1}$; scaling the three-dimensional volume by $h(T)^3$
and controlling its density gives
\[
 \operatorname{Vol}_{g_T^*}(R_{T,i})
 \geq c_n h(T)^3/\delta(T),\qquad
 \operatorname{Vol}_{g(T)}(C_{T,i})\leq C_{\rm cap}h(T)^3.
\]
Here $c_n>0$ is uniform, and $C_{\rm cap}>0$ depends on the selected
standard cap. Let $d_{\rm cap}(K)>0$ and $d_{\rm neck}>0$ be the
applicability thresholds for these two volume estimates. Choose
\[
 d=\min\{d_{\rm cap}(K),d_{\rm neck},c_n/(2C_{\rm cap})\}>0,
 \qquad 0<\overline\delta\leq d,
\]
and require $\delta(T)\leq\overline\delta$. Subtracting the cap bound leaves
$c h(T)^3/\delta(T)$ with $c=c_n/2$ for every cap.
The positive half-necks are pairwise disjoint and disjoint from the retained
region in the terminal carrier.  On the post-surgery carrier the retained
region and the caps cover the slice.  Add the per-cap inequalities and use
the retained-volume identity in \eqref{eq:global-event-volume}; this gives
the stronger multiplicity bound
\begin{equation}
 v_F(T)+n_T c\,h(T)^3/\delta(T)
 \leq\operatorname{Vol}_{g_T^*}(M_T^*)\leq L_T.
 \label{eq:global-all-cap-loss}
\end{equation}
The determinant evolution $\partial_t d\mu=-R\,d\mu$ and the scalar lower
bound $R\geq-6$ give, on a regular slab,

\begin{equation}
 v_F(b)\leq e^{6(b-a)}v_F(a),\qquad a\leq b,
 \label{eq:global-volume-growth}
\end{equation}

and the left-limit comparison concatenates this inequality across events.
Thus $e^{-6t}v_F(t)$ is nonincreasing and pays the weighted sum of all
per-cap losses.  The scale $h^3/\delta$ follows from actual chart volume;
the quadratic expression printed in \cite[Lemma 17.12, pp.~410--411]{MT2007}
is corrected here.

\section{Finite horizons and no finite accumulation}
\label{sec:global-finiteness}

Fix $B>0$, an initial volume bound $V_0<\infty$, and a lower height
$h_{\min}>0$.  On an observed interval contained in $[0,B]$, suppose the
pinching, terminal-strong-neck, disappearing-boundary, and zero-cap controls
hold, while every event satisfies $\delta(T)\leq\bar\delta$ and
$h_{\min}\leq h(T)$.  Set $\sigma=c h_{\min}^3/\bar\delta>0$.
The multiplicity estimate \eqref{eq:global-all-cap-loss} gives, for every
finite set $S$ of observed event times,
\[
 e^{-6B}\sigma\sum_{T\in S}n_T
 \leq\sum_{T\in S} e^{-6T}\sigma n_T\leq V_0.
\]
Choose an integer $N_c\geq e^{6B}V_0/\sigma$ before choosing the flow.
It bounds total cap multiplicity, not just the number of times with a cap.

Let $c(t)$ denote the number of connected components of the slice, and
let $d_T=1$ at a zero-cap or vanishing event and $d_T=0$ otherwise.
Ordinary slab identifications preserve $c(t)$.  Cutting along $n_T$ spheres
can increase the component number by at most $n_T$; filling their boundary
spheres with balls creates no additional component.  At a zero-cap event
the explicit discarded-component witness removes at least one component;
at a vanishing event the finite nonempty incoming slice is replaced by the
empty slice.  Hence at each event
\[
 c(T)+d_T\leq c(T^-)+n_T.
\]
Sum this inequality over the complete event prefix up to an included time
$b$. This prefix is finite by the defining local finiteness of a raw flow
on included compact intervals; the new bound is uniform and also controls
approach to an excluded horizon. The ordinary intervals cancel, giving
$\sum d_T\leq c(0)+\sum n_T$.  Every actual event has
$n_T+d_T\geq1$, so the complete prefix has at most $c(0)+2\sum n_T$
events.  For a requested finite set that omits intermediate events, first
enlarge it to the complete prefix; applying component balance to the
requested set alone would be unjustified.

Finally each initial component contains a unit ball, which stays in that
component because finite Riemannian distance is realized by paths.
Normalization gives volume at least $\omega_3/2$ for such a ball.
The disjoint components therefore satisfy $c(0)\omega_3/2\leq V_0$.
Choose an integer $N_0\geq2V_0/\omega_3$ and put $N=N_0+2N_c$.
This single integer bounds all event sets in every flow with the prescribed
data.  It uses no positive lower volume loss at deletion-only events.

The global restart uses this finite-subset bound through an explicit
$N+1$ contradiction.  Fix an epoch and let $N$ be its uniform bound.  Start
with a stage whose horizon is still below the epoch endpoint. Continuation advances
the same changing-carrier flow to a new stage; if the endpoint has not yet
been reached, the old horizon is a surgery time of the new stage and no later
event occurs in the open gap.  It is therefore fresh in the existing finite
ledger.  Inductively, after $k$ unsuccessful restarts the ledger contains
$k$ distinct observed surgery times.  After $N+1$ such restarts it would be a
finite subset of cardinality $N+1$, contradicting the uniform event bound. Hence one
restart reaches the epoch boundary, with the old terminal slab retained when
the boundary itself is the observed endpoint.  The fixed bound exists for
all restart stages because their actual extension maps preserve initial
volume and the total parameter record: the positive height at the epoch's
right boundary is a common lower height throughout that epoch.
This is the finite continuation argument of
\cite[Section 17.2 and Lemma 17.12, pp.~408--411]{MT2007}.

For a supplied loss certificate on one flow, choose a compact time set
$K\subset\mathbb R$.  Enclose it in a finite interval $[0,H]$ after
intersecting the nonnegative time domain.  The uniform finite-subset bound
implies $(J_F\cap K)$ is finite: an infinite set would contain finite
subsets of arbitrarily large cardinality.  Taking a compact closed interval
around any $T$ yields a radius $d>0$ for which

\[
             J_F\cap(T-d,T+d)\quad\text{is finite}.
\]

Thus the very flow for which the loss estimate was established has locally
finite event times and no finite accumulation. This uses the same volume
loss estimate as the finite-restart argument.

\section{Compatible finite prefixes and the global assembly}
\label{sec:global-assembly}

The global construction is an induction over completed epochs.  A finite
stage contains a flow, its stage observation, the finite parameter prefix, and
comparison maps identifying every earlier stage with a common stage.  When two
stages both contain a time $t$, their comparison maps agree after passing to a
common later stage.  Thus the global identification is independent of the
chosen representative.  The event ledger transports the old event reference
time, retained pre- and post-regions, retention map, and (for an empty slice)
the vanishing-event reference \cite[Definition 15.8, pp.~361--362]{MT2007}.

At stage $n$, apply epoch extension with $p_n$, its selected noncollapsing
and canonical extensions, and the current observation. The finite restart
argument reaches the next epoch boundary. Each new stage is an actual
extension of its predecessor: its comparison maps pull back the metric
exactly on ordinary overlaps and preserve earlier event data. Equation
\eqref{eq:global-profile} also preserves their numerical controls.
The completed epoch horizons tend to infinity. Taking the union of their
representatives under these compatible identifications therefore gives
one raw flow $G$ on $[0,\infty)$, with a well-defined metric
on each ordinary slab and a well-defined event family.  The exact finite
prefix witness at index $i$ is retained for every $i>0$.

There are two applications. Given a finite controlled prefix $F$ on $[0,T)$, with

\begin{align*}
 \mathcal J_F&=[0,T),\\
 \epsilon_F&=\epsilon,\quad C_F=C,\\
 \delta_F(t)&=\delta(t),
 \end{align*}

and the old event policy, pinching, canonical and noncollapsing controls, it
returns an extension of that same $F$ to $[0,\infty)$.  The extension keeps
the schedule agreement, terminal policy, admissibility, compact-interval
finiteness, and the no-finite-accumulation property.  The proof first checks
whether the given flow already reaches the required horizon; otherwise it
forms a completed finite stage chain and applies the controlled extension
operation to the actual branch.  It never chooses a new solution in place of
the given one.

For the normalized branch, let $g_0$ be a normalized initial metric on a compact
nonempty three-manifold $M$ with no embedded projective plane of trivial normal
bundle.  Let $\delta$ be antitone on $[0,\infty)$, positive there, and satisfy

\begin{equation}
 \delta(t)\leq\Delta(j)\quad(t\in E_j).
 \label{eq:global-control-cutoff}
\end{equation}

The normalized initial-flow estimate supplies the first stage, including
the time-zero identification. The
completed chain produces a flow $G$ for the already selected schedule $S$, with

\[
 \mathcal J_G=[0,\infty),\qquad
 \delta_G=\delta,\qquad \epsilon_G=\epsilon,
 \quad C_G=C.
\]

There is a diffeomorphism

\[
 \iota:M\longrightarrow M_0^G
\]

such that for all $x,v,w$

\begin{equation}
 g_G(0)_{\iota(x)}(d\iota_xv,d\iota_xw)=g_0(x)(v,w).
 \label{eq:global-initial-identification}
\end{equation}

The initial carrier is therefore the chosen normalized manifold, up to the
specified smooth identification; later carriers are the actual surgery
carriers.

Apply the volume argument to this same flow, using its admissibility,
pinching, event pre-intervals and discarded components at zero-cap events.
The resulting loss estimate implies finite
volume on compact time sets, the left-limit loss comparison, the cubic scale,
the component-event bound, local finiteness, and no finite accumulation.
Once a slice is empty, every later slice
is empty and no later surgery occurs.

We have proved the following global existence statement, with the same
schedule and control function throughout:

\begin{quote}
for every $\varepsilon_{\rm bound}>0$ there are $K$ and a global schedule
$S$ with $2\epsilon\leq\varepsilon_{\rm bound}$; every admissible
finite controlled prefix satisfying the schedule extends on its own flow;
and every normalized initial metric with the projective-plane exclusion and every
positive antitone $\delta$ satisfying \eqref{eq:global-control-cutoff}
produces a complete global certificate identified with $S$ and $\delta$.
\end{quote}

The resulting flow has canonical neighborhoods, noncollapsing,
pinching, admissibility, compact volume bounds, event-loss data, local
finiteness, and empty-slice permanence.  The source quantity is exactly
$\rho=\delta r$ in the height selector. Extinction will be proved using
the width comparison.

\section{Choosing the flow for the width comparisons}
\label{sec:global-comparison-calibration}

The global existence theorem permits a choice of accuracy bound and of
positive antitone cutoff.  We now make those choices before constructing
the flow, so that its own local topology and surgery comparison maps
satisfy the hypotheses of the extinction argument.

Let $\varepsilon_{\rm top},\varepsilon_{\rm map},
\varepsilon_{\rm hom}>0$ be the thresholds supplied respectively by
local reconstruction, the comparison map and its homotopy theorem.  Put
$e_* =\min\{\varepsilon_{\rm top},\varepsilon_{\rm map},
\varepsilon_{\rm hom}\}$.  With the fixed accuracy factor
$A=10^{14}$, request the global schedule using the bound $2e_*/A$.
The global theorem returns $2\varepsilon\leq2e_*/A$, hence
$A\varepsilon\leq e_*$.  The constants and schedule are chosen before
the initial manifold.  The factor $A$ is the actual uniform loss retained
by the construction; replacing it by $2$ would weaken the hypotheses of
the subsequent local theorems.

Write $\delta_0,R_0>0$ for the metric-surgery cutoffs, and $\Delta(j)$
for the positive antitone schedule, with $\Delta(j)\leq\delta_0$.
Define, before selecting the flow,
\begin{equation}
 \delta_*(t)=\min\left\{
   \frac{\Delta(\lceil32t\rceil_+)}2,
   1,\frac{R_0^{-1/2}}2\right\},
 \label{eq:global-strict-control}
\end{equation}
where $\lceil s\rceil_+$ is the least nonnegative integer at least $s$.
This function is positive and antitone: the ceiling index is nondecreasing,
$\Delta$ is antitone, and taking a minimum preserves this order.
If $t$ belongs to epoch $E_j$, then $j\leq\lceil32t\rceil_+$.
For $j=0$ this is automatic; for $j=n+1$, the lower epoch bound gives
$2^n\leq32t$, and $n+1\leq2^n$.  Therefore
$\delta_*(t)\leq\Delta(j)/2\leq\Delta(j)$, as required for the
normalized global-flow construction.

Apply that construction with this exact control.  The resulting flow has
$\delta(t)=\delta_*(t)$ on its nonnegative domain, so
\[
 \delta(t)\leq\frac{\Delta(0)}2<\Delta(0),\qquad
 \delta(t)\leq\frac{\delta_0}2,\qquad
 \delta(t)\leq1,\qquad
 \delta(t)\leq\frac{R_0^{-1/2}}2.
\]
Its height bound and $r(t)\leq1$ give
\[
 h(t)\leq\delta(t)^2r(t)\leq\delta(t),
 \qquad h(t)\leq\frac{R_0^{-1/2}}2<R_0^{-1/2}.
\]
In particular $\delta(t)<\delta_0$ and $h(t)<R_0^{-1/2}$ hold at
every surgery time, with explicit margins.

The flow retains the chosen schedule's accuracy.  Apply local
reconstruction to its admissibility and the bound
$A\varepsilon\leq\varepsilon_{\rm top}$; apply the comparison-map
theorem at $\varepsilon_{\rm map}$, and then apply the homotopy theorem
to that very comparison map at $\varepsilon_{\rm hom}$.  All three
outputs concern the same flow.  This supplies the local event geometry
and strict comparison bounds needed downstream, without changing the
flow after a target time has been selected.  The choices implement
\cite[Section 15.1.1, pp.~354--355; Definition 15.7 and
Proposition 15.12, pp.~360, 365]{MT2007}; no optimality is claimed for
the explicit constants.

\section{The flow used in extinction}
\label{sec:global-interfaces}

The next chapter consumes the global certificate through the single flow $G$,
its event ledger, and the permanent-empty property.  It may use the
canonical/noncollapsing controls and the exact schedule agreement on each
half-open epoch.  It must not replace $G$ by a different continuation or
identify incoming and outgoing carriers without the event maps.  The
extinction chapter supplies the scalar comparison and the terminal empty
slice; this chapter supplies neither.

The noncollapsing step supplied the stable source and uniform action and
image-volume bounds required by Theorem~\ref{thm:m15-noncollapse}.
The canonical induction used the ancient alternatives and the controlled
backward-limit construction of Section~\ref{sec:singular-blowup}.
At each terminal time, Theorem~\ref{thm:branch-continuation} supplies the
event with the same regular region and height, while
Theorem~\ref{thm:cap-persistence} controls its new caps. The finite-restart
and global-union arguments above preserve those choices and all previous
events. Thus the width comparison can be applied to one fixed flow for
arbitrarily large times.
